%% 
%% Usage instructions
%% 
%% This template can be used in two different ways
%% 1) Use with nowadays pdflatex:  pdflatex diss
%% 2) Use the old way with latex:  latex diss; dvips diss
%%    This way you can use things like psfrag
%% 
%% Bibliography: We use biber as backend. So a complete compilation run looks like this
%% 1) Using pdflatex: pdflatex diss; biber diss; pdflatex diss; pdflatex diss
%% 2) Using Latex: latex diss; biber diss; latex diss; latex diss; dvips diss; ps2pdf14 -dUseFlateCompression#true -dPDFSETTINGS#/prepress -dEmbedAllFonts#true -sPAPERSIZE#a4 diss.ps diss.pdf
%% 
%% 
%% 
%% 
%% 
%% 
%% 
%% 
%% 
%% 
%% If you have any questions, problems or feature request: frank.deinzer@fhws.de
%% 



\documentclass[american,twoside,12pt]{report}

% uncomment if you want to use the debugging features
%\def\DEBUG{debug}

\ifdefined\DEBUG%
	\PassOptionsToPackage{showframe}{geometry}	% DEBUG, FRAME AROUND DOCUMENT
\fi

\usepackage[a4paper,vmargin=35mm,hmargin=25mm]{geometry}
\usepackage{ifpdf}    % check if we are running pdflatex or something else (may be usefull for psfrag things)
\usepackage{setspace}	% baselinestretch?!
\usepackage{times}		% font, times package is a little outdated, but all the new one create problems  
\onehalfspacing
\let\oldChapter=\chapter
\def\chapter{\cleardoublepage\oldChapter}
\pagestyle{headings}
% better monospaced font
\usepackage{inconsolata}

\usepackage[utf8]{inputenc}
\usepackage[T1]{fontenc}

\PassOptionsToPackage{hyphens}{url}
\ifdefined\DEBUG%
	\usepackage[colorlinks]{hyperref}		% DEBUG, COLOR LINKS
	\newcommand{\debugRed}{\color{red}}
	\newcommand{\debugBlue}{\color{blue}}
	\newcommand{\debugBlueB}{\textcolor{blue}}
	\newcommand{\markEq}{ \color{red} }
	\newcommand{\markText}[1]{ {\color{red} #1} }
\else
	\usepackage[hidelinks]{hyperref}		% PRINT
	\newcommand{\debugRed}{}			
	\newcommand{\debugBlue}{}
	\newcommand{\debugBlueB}{}
	\newcommand{\markEq}{}
	\newcommand{\markText}[1]{#1}
\fi%

\usepackage{graphicx}
\usepackage{wrapfig}	% floating wrapped figures
\graphicspath{{images/}}   % graphics search paths, you can give multiple paths
\usepackage{psfrag}

\usepackage{sidecap}	% side-captions
\usepackage[font=small]{caption}	
\usepackage[font=small]{subcaption}	

% math
\usepackage{amsmath}
\usepackage{amssymb}
\usepackage{braket}
\usepackage{mathtools}
% math defines. Add your own defines 
\def\mat#1{\ensuremath{\boldsymbol#1}}
\def\vec#1{\ensuremath{\boldsymbol#1}}
\makeatletter 
\def\argmax{\mathop{\operator@font argmax}}
\def\argmin{\mathop{\operator@font argmin}}
\makeatother


% rotated table
\usepackage{rotating}

% toprule and midrule instead of hline
\usepackage{booktabs}

% https://tex.stackexchange.com/questions/247531/how-to-use-boondox-calligraphic-font-in-latex-without-replacing-mathcal-command
% keep original mathcal, add other fonts using other commands
\DeclareFontFamily{U}{BOONDOX-calo}{\skewchar\font=45 }
\DeclareFontShape{U}{BOONDOX-calo}{m}{n}{
	<-> s*[1.05] BOONDOX-r-calo}{}
\DeclareFontShape{U}{BOONDOX-calo}{b}{n}{
	<-> s*[1.05] BOONDOX-b-calo}{}
\DeclareMathAlphabet{\mathcalboondox}{U}{BOONDOX-calo}{m}{n}
\SetMathAlphabet{\mathcalboondox}{bold}{U}{BOONDOX-calo}{b}{n}
\DeclareMathAlphabet{\mathbcalboondox}{U}{BOONDOX-calo}{b}{n}

% algorithm
\usepackage[plain]{algorithm}
\usepackage{algpseudocode}

% tables, colored table cells
\usepackage{multirow,tabularx}

% sources (bibtex) 
% options: alphabetic, ext-alphabetic, nature, science, ieee
\usepackage[backend=biber,maxnames=99,language=english,style=alphabetic,sorting=nty,sortcites=true,doi=false,isbn=false,backref=true,backrefstyle=two,]{biblatex}
%\DeclareFieldFormat{issuedate}{#1}		% when using ext-alphabetic, remove the () around the year for @article...
\addbibresource{diss.bib}   % add your own bib ressources with more addbibresource

% index
%\usepackage{makeidx}
\usepackage{imakeidx}
\makeindex{}

\hbadness=10000
\clubpenalty = 10000 % schliesst Schusterjungen aus
\widowpenalty = 10000 % schliesst Hurenkinder aus

\usepackage{nicefrac}	% for horizontal fractions (texlive: units package)

\begin{document}

	\cleardoublepage
	
	{
		\thispagestyle{empty}
		\newcommand{\myTitle}{Proof on P=NP}
\newcommand{\myName}{Pax Nuptermann}
\title{\myTitle}

%\date{Lübeck -- 2020}

\author{\myName}

\begin{titlepage}%
	%
	%
	% LOGO
	%
	\noindent
\ifpdf	
	\includegraphics[trim=25 18 25 18, clip]{Logo_Uni_Luebeck_BW.pdf}	%	[trim=left bottom right top, clip]
\else
	\includegraphics[trim=25 18 25 18, clip]{Logo_Uni_Luebeck_BW.eps}	%	[trim=left bottom right top, clip]
\fi
	
	\centering
	\vspace{1.0cm}
	
	\large
	
	%
	{\bf From the Institute of Medical Informatics \par}
	{\bf of the University of Lübeck \par}
	{\bf Director: Prof. Dr. rer. nat. habil. Heinz Handels \par}
	
	\vfill
	
	{\bf \Large ``\myTitle'' \par}
	
	\vspace{0.5cm}
	
	{Dissertation \par}
	{for Fulfillment of \par}
	{Requirements \par}
	{for the Doctoral Degree \par}
	{of the University of Lübeck \par}
		
	\vspace{0.5cm}
		
	{from the Department of Computer Sciences and Technical Engineering \par}
	
	\vfill
	
	{Submitted by \par}
	%
	\vspace{0.5cm}
	%
	{\myName \par}
	{from Turinghausen \par}
	%
	\vspace{0.5cm}
	%
	{Lübeck 202X \par}
	
	
\end{titlepage}


% page 2
\thispagestyle{empty}

	{
		\setlength\parindent{0pt}	% no indent
		\large
		\,\vfill			% empty space needed for vfill
				
		\begin{tabular}{ll}
			First referee:				&	Prof. Dr.-Ing. A. Berta	\\
			Second referee:				&	Prof. Dr.-Ing. B. Ceasar		\\
			Third referee:				&	Prof. Dr.-Ing. C. Dora		\\[8pt]
			Date of oral examination:	&	XX.XX.202X						\\[8pt]
			Approved for printing:		&	Lübeck, XX.XX.202X				\\
		\end{tabular}
		
	}

\newpage
	}

	{
		\thispagestyle{empty}
		\input{chapters/prefaces.tex}
	}
	
	{
		\cleardoublepage
		\pagestyle{plain}
		\pagenumbering{roman} 
		\tableofcontents
		\cleardoublepage
		\pagestyle{empty}
	}
	
	\pagenumbering{arabic}
	\input{chapters/introduction.tex}
	
	\cleardoublepage
	\chapter{Experiments}

\section{Good things}

Lorem ipsum dolor sit amet, consectetuer adipiscing elit. Aenean commodo ligula 
eget dolor. Aenean massa. Cum sociis natoque penatibus et magnis dis parturient 
montes, nascetur ridiculus mus. Donec quam felis, ultricies nec, pellentesque 
eu, pretium quis, sem. Nulla consequat massa quis enim. Donec pede justo, 
fringilla vel, aliquet nec, vulputate eget, arcu. In enim justo, rhoncus ut, 
imperdiet a, venenatis vitae, justo. Nullam dictum felis eu pede mollis pretium. 
Integer tincidunt. Cras dapibus. Vivamus elementum semper nisi. Aenean vulputate 
eleifend tellus. Aenean leo ligula, porttitor eu, consequat vitae, eleifend ac, 
enim. Aliquam lorem ante, dapibus in, viverra quis, feugiat a, tellus. Phasellus 
viverra nulla ut metus varius laoreet. Quisque rutrum. Aenean imperdiet. Etiam 
ultricies nisi vel augue. Curabitur ullamcorper ultricies nisi. Nam eget dui. 
Etiam rhoncus. Maecenas tempus, tellus eget condimentum rhoncus, sem quam semper 
libero, sit amet adipiscing sem neque sed ipsum. Nam quam nunc, blandit vel, 
luctus pulvinar, hendrerit id, lorem. Maecenas nec odio et ante tincidunt 
tempus. Donec vitae sapien ut libero venenatis faucibus. Nullam quis ante. Etiam 
sit amet orci eget eros faucibus tincidunt. Duis leo. Sed fringilla mauris sit 
amet nibh. Donec sodales sagittis magna. Sed consequat, leo eget bibendum 
sodales, augue velit cursus nunc, 

Sed ut perspiciatis unde omnis iste natus error sit voluptatem accusantium 
doloremque laudantium, totam rem aperiam, eaque ipsa quae ab illo inventore 
veritatis et quasi architecto beatae vitae dicta sunt explicabo. Nemo enim ipsam 
voluptatem quia voluptas sit aspernatur aut odit aut fugit, sed quia 
consequuntur magni dolores eos qui ratione voluptatem sequi nesciunt. Neque 
porro quisquam est, qui dolorem ipsum quia dolor sit amet, consectetur, adipisci 
velit, sed quia non numquam eius modi tempora incidunt ut labore et dolore 
magnam aliquam quaerat voluptatem. Ut enim ad minima veniam, quis nostrum 
exercitationem ullam corporis suscipit laboriosam, nisi ut aliquid ex ea commodi 
consequatur? Quis autem vel eum iure reprehenderit qui in ea voluptate velit 
esse quam nihil molestiae consequatur, vel illum qui dolorem eum fugiat quo 
voluptas nulla pariatur? At vero eos et accusamus et iusto odio dignissimos 
ducimus qui blanditiis praesentium voluptatum deleniti atque corrupti quos 
dolores et quas molestias excepturi sint occaecati cupiditate non provident, 
similique sunt in culpa qui officia deserunt mollitia animi, id est laborum et 
dolorum fuga. Et harum quidem rerum facilis est et expedita distinctio. Nam 
libero tempore, cum soluta nobis est eligendi optio cumque nihil impedit quo 
minus id quod maxime placeat facere 

\section{Better things}

Li Europan lingues es membres del sam familie. Lor separat existentie es un 
myth. Por scientie, musica, sport etc, litot Europa usa li sam vocabular. Li 
lingues differe solmen in li grammatica, li pronunciation e li plu commun 
vocabules. Omnicos directe al desirabilite de un nov lingua franca: On refusa 
continuar payar custosi traductores. At solmen va esser necessi far uniform 
grammatica, pronunciation e plu sommun paroles. Ma quande lingues coalesce, li 
grammatica del resultant lingue es plu simplic e regulari quam ti del coalescent 
lingues. Li nov lingua franca va esser plu simplic e regulari quam li existent 
Europan lingues. It va esser tam simplic quam Occidental in fact, it va esser 
Occidental. A un Angleso it va semblar un simplificat Angles, quam un skeptic 
Cambridge amico dit me que Occidental es. Li Europan lingues es membres del sam 
familie. Lor separat existentie es un myth. Por scientie, musica, sport etc, 
litot Europa usa li sam vocabular. Li lingues differe solmen in li grammatica, 
li pronunciation e li plu commun vocabules. Omnicos directe al desirabilite de 
un nov lingua franca: On refusa continuar payar custosi traductores. At solmen 
va esser necessi far uniform grammatica, pronunciation e plu sommun paroles. 

\begin{figure}
 \begin{center}{
 \ifpdf
 		\includegraphics[width=0.7\textwidth]{PNG_example}
\else
    \psfrag{s0}{\footnotesize $\tau$}
    \psfrag{s1}{\footnotesize $Q( {\vec{s}^1} , {\vec{a}^1} )$}
    \psfrag{s2}{\footnotesize $Q( {\vec{s}^2} , {\vec{a}^2} )$}
    \psfrag{s3}{\footnotesize $Q( {\vec{s}^3} , {\vec{a}^3} )$}
    \psfrag{s4}{\footnotesize $Q( {\vec{s}^4} , {\vec{a}^4} )$}
    \psfrag{s5}{\footnotesize $Q( {\vec{s}^5} , {\vec{a}^5} )$}
    \psfrag{s6}{\footnotesize $Q( {\vec{s}^6} , {\vec{a}^6} )$}
    \psfrag{s7}{\footnotesize $Q( {\vec{s}} , {\vec{a}} )$}
    \psfrag{s8}{\footnotesize $\Phi$}
    \psfrag{s9}{\footnotesize $Q( {\vec{s}^7} , {\vec{a}^7} )$}
    \psfrag{w}{\small $W$}
		\includegraphics[width=0.7\textwidth]{EPS_example}
\fi
  }\end{center}
\ifpdf
	\caption[Exmaple image: Either png or eps]{This is only a png image. If using latex and dvips one could do more sophisticated things with text replacement. See sources for example} 
\else
	\caption[Exmaple image: Either png or eps]{This is a fancy image: eps format with label replacement with psfrag. This way you get correct fonts and math into images. But does not work with pdflatex. Only with latex and dvips.}
\fi
\end{figure}

Lorem ipsum dolor sit amet, consectetuer adipiscing elit. Aenean commodo ligula 
eget dolor. Aenean massa. Cum sociis natoque penatibus et magnis dis parturient 
montes, nascetur ridiculus mus. Donec quam felis, ultricies nec, pellentesque 
eu, pretium quis, sem. Nulla consequat massa quis enim. Donec pede justo, 
fringilla vel, aliquet nec, vulputate eget, arcu. In enim justo, rhoncus ut, 
imperdiet a, venenatis vitae, justo. Nullam dictum felis eu pede mollis pretium. 
Integer tincidunt. Cras dapibus. Vivamus elementum semper nisi. Aenean vulputate 
eleifend tellus. Aenean leo ligula, porttitor eu, consequat vitae, eleifend ac, 
enim. Aliquam lorem ante, dapibus in, viverra quis, feugiat a, tellus. Phasellus 
viverra nulla ut metus varius laoreet. Quisque rutrum. Aenean imperdiet. Etiam 
ultricies nisi vel augue. Curabitur ullamcorper ultricies nisi. Nam eget dui. 
Etiam rhoncus. Maecenas tempus, tellus eget condimentum rhoncus, sem quam semper 
libero, sit amet adipiscing sem neque sed ipsum. Nam quam nunc, blandit vel, 
luctus pulvinar, hendrerit id, lorem. Maecenas nec odio et ante tincidunt 
tempus. Donec vitae sapien ut libero venenatis faucibus. Nullam quis ante. Etiam 
sit amet orci eget eros faucibus tincidunt. Duis leo. Sed fringilla mauris sit 
amet nibh. Donec sodales sagittis magna. Sed consequat, leo eget bibendum 
sodales, augue velit cursus nunc, 

\section{Best things}

\subsection{Best part1}

\paragraph{A paragraph} Sed ut perspiciatis unde omnis iste natus error sit voluptatem accusantium 
doloremque laudantium, totam rem aperiam, eaque ipsa quae ab illo inventore 
veritatis et quasi architecto beatae vitae dicta sunt explicabo. Nemo enim ipsam 
voluptatem quia voluptas sit aspernatur aut odit aut fugit, sed quia 
consequuntur magni dolores eos qui ratione voluptatem sequi nesciunt. Neque 
porro quisquam est, qui dolorem ipsum quia dolor sit amet, consectetur, adipisci 
velit, sed quia non numquam eius modi tempora incidunt ut labore et dolore 
magnam aliquam quaerat voluptatem. Ut enim ad minima veniam, quis nostrum 
exercitationem ullam corporis suscipit laboriosam, nisi ut aliquid ex ea commodi 
consequatur? Quis autem vel eum iure reprehenderit qui in ea voluptate velit 
esse quam nihil molestiae consequatur, vel illum qui dolorem eum fugiat quo 
voluptas nulla pariatur? At vero eos et accusamus et iusto odio dignissimos 
ducimus qui blanditiis praesentium voluptatum deleniti atque corrupti quos 
dolores et quas molestias excepturi sint occaecati cupiditate non provident, 
similique sunt in culpa qui officia deserunt mollitia animi, id est laborum et 
dolorum fuga. Et harum quidem rerum facilis est et expedita distinctio. Nam 
libero tempore, cum soluta nobis est eligendi optio cumque nihil impedit quo 
minus id quod maxime placeat facere 

\subsection{Best part2}

Li \index{Europan} lingues es membres del sam familie. Lor separat existentie es un 
myth. Por scientie, musica, sport etc, litot Europa usa li sam vocabular. Li 
lingues differe solmen in li grammatica, li pronunciation e li plu commun 
vocabules. Omnicos directe al desirabilite de un nov lingua franca: On refusa 
continuar payar custosi traductores. At solmen va esser necessi far uniform 
grammatica, pronunciation e plu sommun paroles. Ma quande lingues coalesce, li 
grammatica del resultant lingue es plu simplic e regulari quam ti del coalescent 
lingues. Li nov lingua franca va esser plu simplic e regulari quam li existent 
Europan lingues. It va esser tam simplic quam Occidental in fact, it va esser 
Occidental. A un Angleso it va semblar un simplificat Angles, quam un skeptic 
Cambridge amico dit me que Occidental es. Li Europan lingues es membres del sam 
familie. Lor separat existentie es un myth. Por scientie, musica, sport etc, 
litot Europa usa li sam vocabular. Li lingues differe solmen in li grammatica, 
li pronunciation e li plu commun vocabules. Omnicos directe al desirabilite de 
un nov lingua franca: On refusa continuar payar custosi traductores. At solmen 
va esser necessi far uniform grammatica, pronunciation e plu sommun paroles. 

 

	\cleardoublepage
	\input{chapters/summary.tex}

	\cleardoublepage
	\input{chapters/futureWork.tex}

	% List of figures
	\phantomsection
	\addcontentsline{toc}{chapter}{\listfigurename}
	\listoffigures
	\cleardoublepage

	% List of tables
	\phantomsection
	\addcontentsline{toc}{chapter}{\listtablename}
	\renewcommand{\listtablename}{List of Tables\vspace{-5mm}}	
	\listoftables
	\cleardoublepage

	% Literatur
	\phantomsection
	\addcontentsline{toc}{chapter}{\bibname}
	%\printbibliography[title={All Bib}]
	\printbibliography
	\cleardoublepage
	
	\begin{appendix}
	\cleardoublepage
	\input{chapters/appendix.tex}
	\end{appendix}
	
	\cleardoublepage
	%% index
	\phantomsection
	\printindex
	
	\cleardoublepage

\end{document}
